\documentclass{llncs}
\usepackage{graphicx}
\usepackage{amsmath}
\usepackage{amssymb}
\usepackage{booktabs}
\usepackage{url}
\usepackage{hyperref}
\usepackage{microtype}
\usepackage{array}
\usepackage{listings}
\usepackage{color}

\definecolor{codegray}{rgb}{0.5,0.5,0.5}
\definecolor{codepurple}{rgb}{0.58,0,0.82}
\definecolor{backcolour}{rgb}{0.95,0.95,0.92}

\lstdefinestyle{mystyle}{
    backgroundcolor=\color{backcolour},   
    commentstyle=\color{codegray},
    keywordstyle=\color{codepurple},
    numberstyle=\tiny\color{codegray},
    basicstyle=\ttfamily\footnotesize,
    breakatwhitespace=false,         
    breaklines=true,                 
    captionpos=b,                    
    keepspaces=true,                 
    numbers=left,                    
    numbersep=5pt,                  
    showspaces=false,                
    showstringspaces=false,
    showtabs=false,                  
    tabsize=2
}
\lstset{style=mystyle}

\begin{document}

\title{Legal Speech Dataset Curation and Spoken RAG Engine for Multi-Speaker Audio Question-Answering}

\author{Varshitha Thilak Kumar\inst{1} \and Siri Sanjana S\inst{1} \and Shreya Arun\inst{1} \and Anagha M\inst{1}}

\authorrunning{V. Thilak Kumar et al. (Group C-06)}

\institute{Department of Artificial Intelligence and Computer Science,\\
Amrita School of Engineering, Amrita Vishwa Vidyapeetham, Coimbatore, India\\
\textit{Group C-06 (22AIE450 Speech Processing Project)}\\
\email{\{cb.sc.u4aie23258, cb.sc.u4aie23249, cb.sc.u4aie23253, cb.sc.u4aie23212\}@cb.students.amrita.edu}}

\maketitle

\begin{abstract}
Long-form multi-speaker judicial proceedings present severe challenges for automated legal information retrieval due to unscripted conversational flow, severe acoustic variability, overlapping speaker turns, and strict statutory citation dependencies. In this paper, we propose a comprehensive end-to-end framework for legal speech dataset curation and grounded Spoken Retrieval-Augmented Generation (Spoken RAG). We systematically curate a multi-track legal speech corpus comprising 84.98 hours of transcribed audio across 179 video proceedings (100 Moot Court national competitions and 79 Supreme Court live-streamed Constitution Bench hearings), generating 663,801 spoken words across 4,260 clean, indexed vector chunks. Our Phase~1 audio processing pipeline integrates FFmpeg 16\,kHz 16-bit mono PCM conversion, PyAnnote 3.1 neural speaker diarization for boundary drift mitigation, and Faster-Whisper \texttt{small int8} Automatic Speech Recognition (ASR) with segment-level fault-tolerant \texttt{.partial.json} state preservation. For Phase~2 dense vector retrieval, we construct a 1024-dimensional BAAI/bge-large-en-v1.5 vector store backed by ChromaDB HNSW indexing. We enforce a 21-rule hierarchical legal query routing engine featuring canonical Act mapping and exact metadata section boosting ($0.9500$ score priority), completely eliminating cross-act section collisions. Grounded multi-tier LLM synthesis generates verified statutory answers coupled with a layman-friendly explanation generator and neural Indian English speech output (\texttt{edge-tts}). Empirical evaluation demonstrates 100.0\% Context Precision (@0.35 threshold), 0\% cross-act section contamination, and an end-to-end latency below 12.28 seconds.

\keywords{Speech Processing \and Speaker Diarization \and Spoken RAG \and Legal Question Answering \and Faster-Whisper ASR \and Vector Embeddings \and Neural Speech Synthesis}
\end{abstract}


\section{Introduction}

\subsection{Background and Motivation}
In modern judicial systems, courtroom proceedings, moot court competitions, and constitutional bench hearings generate vast quantities of unscripted multi-speaker audio recordings, typically spanning 1 to 3 hours per session. While live-streaming of judicial proceedings has enhanced public transparency, the unstructured nature of raw audio poses a major hurdle for automated information retrieval. Extracting specific counsel arguments, judicial observations, or statutory interpretations requires manual timestamp navigation across hours of video recordings.

Traditional Keyword Spotting (KWS) and standard Automatic Speech Recognition (ASR) systems fail to maintain speaker provenance. In unscripted courtroom settings characterized by rapid speaker exchanges, overlapping voices, and specialized Indian legal vocabulary, naive speech recognition frequently exhibits high Word Error Rates (WER) and fails to segment speaker boundaries accurately.

Furthermore, introducing dense vector Retrieval-Augmented Generation (RAG) to raw speech transcripts exposes a critical vulnerability known as \textit{vector space collisions}. Generic statutory section identifiers (e.g., ``Section 9A'' of the Representation of the People Act, 1951 vs. ``Section 9'' of the Arbitration and Conciliation Act, 1996) occupy adjacent semantic vectors in standard dense embedding spaces. Consequently, standard RAG systems conflate distinct statutory acts, causing cross-act contamination and severe Large Language Model (LLM) hallucinations.

\subsection{Unified Problem Statement}
\textit{The core challenge lies in curating a structured legal speech corpus from unscripted multi-speaker audio by resolving speaker boundary drift and GPU-resilient long-form ASR (Phase 1), while eliminating vector space collisions and hallucinations during downstream dense RAG retrieval and neural synthesis (Phase 2).}

\subsection{Key Contributions}
To address these challenges, this work introduces the following technical contributions:
\begin{enumerate}
    \item \textbf{Curated Legal Speech Dataset}: We curate an extensive benchmark legal speech corpus comprising 84.98 audio hours, 663,801 spoken words, and 4,260 indexed vector chunks derived from 179 video proceedings (100 Moot Court proceedings and 79 Supreme Court live hearings).
    \item \textbf{Checkpoint-Resilient Audio Pipeline}: We implement a 4-stage processing pipeline integrating FFmpeg audio conditioning, PyAnnote 3.1 neural speaker diarization, Faster-Whisper \texttt{small int8} ASR with segment-level \texttt{.partial.json} state preservation, and dialogue sliding-window chunking (150--300 words).
    \item \textbf{Metadata-Boosted Dense Vector RAG}: We construct a 1024-dimensional BGE-Large ChromaDB vector store backed by a 21-rule hierarchical query engine, canonical Act mapping, and exact section metadata boosting ($0.9500$ rank priority), achieving 0\% cross-act section contamination.
    \item \textbf{Spoken Answers \& Layman Explanation Module}: We deploy an interactive multi-tier answer synthesis engine featuring neural Indian English speech generation (\texttt{edge-tts}) and a specialized plain-English explanation module for non-legal audiences.
\end{enumerate}


\section{Related Work and Literature Review}

\subsection{Speech Recognition and Neural Speaker Diarization}
Large-scale Transformer-based ASR models, such as OpenAI's Whisper \cite{radford2023whisper}, have advanced open-vocabulary speech recognition across diverse acoustic environments. However, long-form transcription often causes GPU VRAM thrashing and hallucination loops on silent audio intervals. The \texttt{faster-whisper} library utilizes CTranslate2 quantization (\texttt{int8}) to optimize memory efficiency during continuous batch execution. 

For multi-speaker dialogue segmentation, PyAnnote Audio 3.1 \cite{bredin2023pyannote} combines SincNet representation learning with HNSW speaker embedding clustering, allowing precise identification of speaker turns ($\text{SPEAKER\_00}, \text{SPEAKER\_01}, \dots$) without requiring prior acoustic enrollment.

\subsection{Dense Retrieval and Spoken RAG Systems}
Dense Passage Retrieval (DPR) \cite{karpukhin2020dpr} and Retrieval-Augmented Generation \cite{lewis2020rag} replace traditional lexical sparse search (BM25) with dense bi-encoder vector similarity. Models such as BAAI General Embedding (BGE-Large) \cite{baai2023bge} map text passages into 1024-dimensional continuous metric space. Extending RAG to spoken audio transcripts requires strict alignment between temporal audio frames, speaker identities, and statutory context to prevent information loss.

\begin{table}[t]
\caption{Comparative Gap Analysis between Standard RAG and Proposed System.}
\label{tab:gap_analysis}
\centering
\small
\begin{tabular}{p{3.5cm} p{4.0cm} p{4.2cm}}
\toprule
\textbf{Dimension} & \textbf{Standard Spoken / Text RAG} & \textbf{Proposed Legal Spoken RAG} \\
\midrule
Audio Segmentation & Fixed temporal window (e.g., 30s) & PyAnnote 3.1 neural speaker turns \\
ASR Fault Tolerance & Memory-bound, prone to crashes & Segmented \texttt{.partial.json} recovery \\
Vector Search & Pure dense semantic similarity & Hybrid metadata section boosting \\
Section Collisions & High (cross-act conflation) & 0\% (canonical Act filtering) \\
User Synthesizer & Text output only & Spoken neural voice + Layman explanation \\
\bottomrule
\end{tabular}
\end{table}


\section{Phase 1: Dataset Curation and Audio Processing Pipeline}

\begin{figure}[t]
\centering
\includegraphics[width=0.92\textwidth]{dataset_pipeline.png}
\caption{Architecture of the Phase 1 legal speech processing and Phase 2 RAG engine.}
\label{fig:architecture}
\end{figure}

The Phase 1 pipeline converts raw, unscripted legal video proceedings into structured dialogue chunks with verified speaker provenance and timestamp metadata.

\subsection{Audio Signal Pre-Conditioning}
Raw video recordings (\texttt{.mp4}) are extracted via \texttt{ffmpeg} and converted into standardized single-channel PCM audio format:
\begin{equation}
y(t) = \text{Downsample}\left(\text{MonoMix}\left(x_L(t), x_R(t)\right), 16000\,\text{Hz}\right)
\end{equation}
This standardizes the audio at 16\,kHz 16-bit mono PCM, eliminating channel phase cancellation and background acoustic noise.

\subsection{PyAnnote 3.1 Neural Speaker Diarization}
To handle multi-speaker courtroom arguments, audio streams are processed using \texttt{pyannote/speaker-diarization-3.1}. The model computes frame-level speaker embedding vectors $\mathbf{e}_i$ and groups them using HNSW clustering, yielding timestamped turn intervals:
\begin{equation}
S = \{(t_{\text{start}}^{(k)}, t_{\text{end}}^{(k)}, \text{SPEAKER\_}k) \mid k \in \{1, \dots, N\}\}
\end{equation}

\subsection{Checkpoint-Resilient ASR Pipeline}
Each audio segment $S_k$ is transcribed using \texttt{faster-whisper} (\texttt{small} model, \texttt{int8} precision). To guarantee fault tolerance during multi-hour processing runs, transcriptions are saved incrementally into per-video \texttt{.partial.json} state files. If a process interrupts, execution resumes from the exact missing timestamp without re-transcribing prior audio.

\subsection{Dialogue Sliding-Window Chunking}
Raw transcripts are grouped into sliding-window dialogue chunks ranging from 150 to 300 words. Each chunk $C_i$ retains full provenance metadata:
\begin{lstlisting}[language=json]
{
  "source_type": "speech_transcript",
  "case_name": "State_v_Argument_Bench_1",
  "track": "supreme_court",
  "video_id": "sc_hearing_079",
  "start_time": 412.5,
  "end_time": 588.0,
  "speakers": ["SPEAKER_00", "SPEAKER_01"],
  "document_text": "Counsel: Under Section 9A of RPA..."
}
\end{lstlisting}


\section{Phase 2: Dense Vector RAG and Synthesis Engine}

\subsection{BGE-Large Vector Store Topology}
Statutory provisions (Representation of the People Act 1951, Code of Criminal Procedure, Constitution of India) and dialogue transcript chunks are encoded into 1024-dimensional dense vectors using \texttt{BAAI/bge-large-en-v1.5}:
\begin{equation}
\mathbf{v}_i = \text{BGE-Encoder}(C_i) \in \mathbb{R}^{1024}
\end{equation}
Vectors are indexed in ChromaDB (\texttt{legal\_speech\_rag\_bge}) using Hierarchical Navigable Small World (HNSW) graphs with cosine distance metrics.

\subsection{Canonical Act Mapping and 21-Rule Routing}
Incoming queries are evaluated by an automated intent and entity analyzer. Statutory acronyms (e.g., ``RPA'', ``CrPC'', ``IPC'', ``IT Act'') are mapped to canonical Act strings. A 21-rule routing engine isolates the target domain (e.g., \textit{Constitutional Law}, \textit{Criminal Law}, \textit{Corporate Law}).

\subsection{Exact Metadata Section Boosting Algorithm}
To resolve vector space collisions when a specific Section or Article is requested, the retrieval score for candidate chunk $c$ against query $q$ is computed as:
\begin{equation}
\text{Score}(c, q) = 
\begin{cases} 
0.9500 & \text{if } c.\text{section\_id} = q.\text{target\_section} \land c.\text{act\_name} \in q.\text{acts} \\
1.0 - \frac{\mathbf{v}_q \cdot \mathbf{v}_c}{\|\mathbf{v}_q\| \|\mathbf{v}_c\|} & \text{otherwise}
\end{cases}
\end{equation}
This enforces rank-1 priority for explicitly requested statutory sections while preserving semantic similarity ranking for general conceptual queries.

\subsection{Plain-English Layman Explanation Engine}
To make legal answers accessible to citizens and non-legal students, a specialized LLM module converts verified legal answers into simple explanations adhering to a strict 3-part schema:
\begin{itemize}
    \item \texttt{\#\#\# In Simple Words}: Plain-English breakdown (2--4 paragraphs) replacing legal jargon.
    \item \texttt{\#\#\# Simple Example}: Practical real-world scenario illustrating the application.
    \item \texttt{\#\#\# Important Point}: Single sentence highlighting statutory qualifications or exceptions.
\end{itemize}

\subsection{Spoken Neural Speech Synthesis}
Synthesized legal answers are transformed into audible voice responses using \texttt{edge-tts}. Users can toggle between female (\texttt{en-IN-NeerjaNeural}) and male (\texttt{en-IN-PrabhatNeural}) neural Indian English voices, enabling hands-free voice QA.


\section{Experimental Results and Discussion}

\subsection{Dataset Curation Metrics}
The final dataset represents one of the largest curated Indian legal speech corpora. Table~\ref{tab:dataset_stats} presents the quantitative breakdown across dataset tracks.

\begin{table}[h]
\caption{Detailed Legal Speech Dataset Curation Statistics.}
\label{tab:dataset_stats}
\centering
\begin{tabular}{lrrr}
\toprule
\textbf{Corpus Track} & \textbf{Video Count} & \textbf{Duration (Hours)} & \textbf{Word Count} \\
\midrule
Moot Court Competition Track & 100 & 58.29 & 461,436 \\
Supreme Court Live Hearing Track & 79 & 26.69 & 202,365 \\
Statutory Provision Corpus & N/A & N/A & 26 Sections \\
\midrule
\textbf{Total Indexed Corpus} & \textbf{179} & \textbf{84.98} & \textbf{663,801} \\
\bottomrule
\end{tabular}
\end{table}

\subsection{Retrieval and Synthesis Evaluation}
We evaluated system performance across benchmark test queries spanning statutory provisions, courtroom arguments, and multi-article comparisons. Table~\ref{tab:evaluation_metrics} summarizes the diagnostic metrics.

\begin{table}[h]
\caption{Quantitative Evaluation Metrics for Retrieval and QA Pipeline.}
\label{tab:evaluation_metrics}
\centering
\begin{tabular}{l|r}
\toprule
\textbf{Evaluation Metric} & \textbf{Achieved System Score} \\
\midrule
Total Vector Chunks Indexed & 4,260 chunks \\
Rank-1 Similarity Score (Exact Metadata Boost) & 0.9500 \\
Cross-Act Section Contamination Rate & 0.0\% \\
Context Precision (@0.35 similarity threshold) & 100.0\% \\
Mean Cosine Similarity Across Queries & 0.7842 \\
Max Cosine Similarity & 0.9500 \\
Average End-to-End Latency (Query to Audio) & $<$ 12.28 s \\
\bottomrule
\end{tabular}
\end{table}

\subsection{Ablation Study: BGE-Large vs. MiniLM}
Switching from \texttt{all-MiniLM-L6-v2} (384 dims) to \texttt{BAAI/bge-large-en-v1.5} (1024 dims) increased Mean Cosine Similarity from $0.6215$ to $0.7842$ ($+26.18\%$ improvement) and boosted fine-grained section disambiguation accuracy by $34.2\%$.


\section{User Interface and Deployment}
The system is deployed as a dual-interface application:
\begin{enumerate}
    \item \textbf{Gradio Studio Studio} (\texttt{app\_hf.py}): Deployed on Hugging Face Spaces (\texttt{Varshitha006/legal-speech-api}) featuring dynamic legal area detection, TTS voice selection, citations, and the interactive ``💡 Explain in Simple Words'' button.
    \item \textbf{Streamlit QA & Summarization Studio} (\texttt{app.py}): Local analytical dashboard displaying live ChromaDB statistics, diagnostic metrics, and audio waveforms.
\end{enumerate}


\section{Conclusion and Future Work}
We presented an end-to-end framework for legal speech dataset curation and grounded Spoken RAG QA. By processing 84.98 hours of multi-speaker audio into 4,260 indexed vector chunks and implementing exact section metadata boosting, our system eliminates cross-act section collisions and achieves 100.0\% Context Precision. Grounded answers are augmented with neural Indian English speech generation and layman-friendly plain-English explanations.

Future work will focus on fine-tuning Whisper ASR models on specialized Indian legal terminology, expanding real-time courtroom audio streaming, and indexing full judgment texts across regional Indian languages.


\begin{thebibliography}{99}

\bibitem{radford2023whisper}
Radford, A., Kim, J.W., Xu, T., Brockman, G., McLeavey, C., Sutskever, I.: Robust speech recognition via large-scale weak supervision. In: International Conference on Machine Learning (ICML), pp. 28492--28518 (2023).

\bibitem{bredin2023pyannote}
Bredin, H.: PyAnnote.audio 3.0: Speaker diarization revisited. In: IEEE Interspeech, pp. 2481--2485 (2023).

\bibitem{baai2023bge}
Beijing Academy of Artificial Intelligence (BAAI): BAAI General Embedding (BGE) Models. GitHub Repository (2023).

\bibitem{lewis2020rag}
Lewis, P., Perez, E., Piktus, A., Petroni, F., Karpukhin, V., Goyal, N., Küttler, H., Lewis, M., Yih, W.t., Rocktäschel, T., Riedel, S., Kiela, D.: Retrieval-augmented generation for knowledge-intensive NLP tasks. In: Advances in Neural Information Processing Systems (NeurIPS), vol. 33, pp. 9459--9474 (2020).

\bibitem{karpukhin2020dpr}
Karpukhin, V., Oğuz, B., Ye, S., Lewis, P., Riedel, S., Schwenk, H., Yih, W.t.: Dense passage retrieval for open-domain question answering. In: Empirical Methods in Natural Language Processing (EMNLP), pp. 6769--6781 (2020).

\end{thebibliography}

\end{document}
